% Options for packages loaded elsewhere
\PassOptionsToPackage{unicode}{hyperref}
\PassOptionsToPackage{hyphens}{url}
%
\documentclass[letterpaper, 10pt, journal, twoside]{IEEEtran}
\usepackage{amsmath,amssymb}
\usepackage{iftex}
\ifPDFTeX
  \usepackage[T1]{fontenc}
  \usepackage[utf8]{inputenc}
  \usepackage{textcomp} % provide euro and other symbols
\else % if luatex or xetex
  \usepackage{unicode-math} % this also loads fontspec
  \defaultfontfeatures{Scale=MatchLowercase}
  \defaultfontfeatures[\rmfamily]{Ligatures=TeX,Scale=1}
\fi
\usepackage{lmodern}
\ifPDFTeX\else
  % xetex/luatex font selection
\fi
% Use upquote if available, for straight quotes in verbatim environments
\IfFileExists{upquote.sty}{\usepackage{upquote}}{}
\IfFileExists{microtype.sty}{% use microtype if available
  \usepackage[]{microtype}
  \UseMicrotypeSet[protrusion]{basicmath} % disable protrusion for tt fonts
}{}
\makeatletter
\@ifundefined{KOMAClassName}{% if non-KOMA class
  \IfFileExists{parskip.sty}{%
    \usepackage{parskip}
  }{% else
}{% if KOMA class
  \KOMAoptions{parskip=half}}
\makeatother
\usepackage{xcolor}
\usepackage{longtable,booktabs,array}
\usepackage{calc} % for calculating minipage widths
% Correct order of tables after \paragraph or \subparagraph
\usepackage{etoolbox}
\makeatletter
\patchcmd\longtable{\par}{\if@noskipsec\mbox{}\fi\par}{}{}
\makeatother
% Allow footnotes in longtable head/foot
\IfFileExists{footnotehyper.sty}{\usepackage{footnotehyper}}{\usepackage{footnote}}
\makesavenoteenv{longtable}
\usepackage{graphicx}
\makeatletter
\def\maxwidth{\ifdim\Gin@nat@width>\linewidth\linewidth\else\Gin@nat@width\fi}
\def\maxheight{\ifdim\Gin@nat@height>\textheight\textheight\else\Gin@nat@height\fi}
\makeatother
% Scale images if necessary, so that they will not overflow the page
% margins by default, and it is still possible to overwrite the defaults
% using explicit options in \includegraphics[width, height, ...]{}
\setkeys{Gin}{width=\maxwidth,height=\maxheight,keepaspectratio}
% Set default figure placement to htbp
\makeatletter
\def\fps@figure{htbp}
\makeatother
\ifLuaTeX
  \usepackage{luacolor}
  \usepackage[soul]{lua-ul}
\else
  \usepackage{soul}
\fi
\setlength{\emergencystretch}{3em} % prevent overfull lines
\providecommand{\tightlist}{%
  \setlength{\itemsep}{0pt}\setlength{\parskip}{0pt}}
\ifLuaTeX
  \usepackage{selnolig}  % disable illegal ligatures
\fi
\IfFileExists{bookmark.sty}{\usepackage{bookmark}}{\usepackage{hyperref}}
\IfFileExists{xurl.sty}{\usepackage{xurl}}{} % add URL line breaks if available
\urlstyle{same}
\hypersetup{
  hidelinks,
  pdfcreator={LaTeX via pandoc}}

\author{}
\date{}

\begin{document}

\textbf{Continuous-Time Variational Latent Force Modeling}

\textbf{for Real-Time Adaptive Control of Soft and Deformable Robots}

\textbf{Rainer Rodrigues}

\emph{Independent Researcher}

rainerrodrigues16@gmail.com

Code:
\href{https://github.com/rainerrodrigues/VLFM_SoftRobotics}{\ul{github.com/rainerrodrigues/VLFM\_SoftRobotics}}

\textbf{Abstract}

\begin{quote}
Soft and deformable robots hold great potential for safe, human-centric
applications, but their infinite-degree-of-freedom dynamics and
nonlinear material properties defy traditional analytical modelling.
Nonparametric methods such as Gaussian Process (GP) regression and
Bayesian system identification offer principled uncertainty
quantification, yet their O(N³) computational complexity and reliance on
offline batch processing preclude real-time embedded control. Recent
active exploration methods (e.g. SoftAE) operate episodically and lack
tight integration with online, safety-critical Model Predictive Control
(MPC). This work introduces a real-time, nonparametric, structure-aware
adaptation framework for deformable systems. Sparse Recursive GP
State-Space Models (RGPSSM) process streaming sensor data and update
hyperparameters and inducing points on the fly; variable stiffness and
morphological shifts are captured through latent state inference rather
than discarded as residual noise. A Variational Latent Force Model
(VLFM) implemented in Julia with Turing.jl and the SciML ecosystem
achieves ∼ 480\,Hz inference with a trajectory-tracking MSE of 0.0107
which is a 92\,\% reduction over a sparse GP baseline and 31\,\% over a
Neural ODE baseline --- confirming viability for real-time, embedded
soft-robot control.

\textbf{Index Terms --- }\emph{Gaussian processes, Bayesian inference,
soft robotics, latent force models, model predictive control,
differentiable programming, Julia, SciML, real-time learning.}
\end{quote}

\textbf{I. INTRODUCTION}

Autonomous robots operating alongside humans expose a fundamental
limitation of classical rigid-body systems: they lack the compliance
required for safe physical interaction. Soft robots, whose bodies
comprise highly deformable continuum structures, represent a
transformative paradigm suited to medical, rehabilitation, and wearable
applications. Yet this compliance introduces a formidable control
challenge: governing dynamics exhibit infinite degrees of freedom,
viscoelastic nonlinearities, and hyperelastic material behaviour that
resist accurate first-principles modelling~{[}1{]},~{[}2{]}.

Data-driven machine learning offers a principled alternative, but
deploying such models in closed-loop robotic systems entails three
critical algorithmic bottlenecks. This paper addresses all three through
a unified, real-time, nonparametric, structure-aware adaptation
framework.

\emph{\textbf{A. Algorithmic Challenges for Real-Time Learning}}

\emph{1) Computational Bottleneck}

Exact GP regression scales as O(N³), precluding real-time
control~{[}3{]},~{[}4{]}. Sparse Recursive GP State-Space
Models~{[}5{]}--{[}7{]} overcome this by dynamically managing a fixed
budget of inducing points on streaming data, preventing catastrophic
forgetting on resource-constrained embedded hardware~{[}8{]},~{[}9{]}.

\emph{2) Non-Stationary Morphological Shifts}

Soft robots undergo continuous changes in stiffness, geometry, and
damping that invalidate fixed-model assumptions~{[}1{]}. Rather than
treating these as residual noise, the proposed framework performs online
hyperparameter optimisation and latent state inference directly from
sensor streams~{[}5{]},~{[}10{]}, enabling adaptation to novel operating
domains beyond the reach of episodic batch-training~{[}11{]},~{[}12{]}.

\emph{3) Uncertainty-Aware Predictive Control}

Propagating the GP posterior through a full MPC receding horizon in
real-time requires reformulating the stochastic optimisation
problem~{[}13{]},~{[}14{]}. This is achieved via Model Predictive Path
Integral (MPPI) sampling~{[}15{]} and chance-constraint reformulation
into convex Second-Order Cone Programs (SOCPs)~{[}13{]}.

\emph{\textbf{B. Why Real-Time Learning is Essential}}

Soft robots are uniquely suited to safety-critical biomedical tasks:
wearable robotics eliminates the stiffness and misalignment of rigid
exoskeletons~{[}16{]}, and surgical soft robots navigate confined
anatomical spaces without harming tissue~{[}16{]},~{[}17{]}. Their
compliance enables morphological computation absorbing interaction
complexity in body mechanics thereby reducing the
controller\textquotesingle s planning burden~{[}2{]}. Because
first-principles models fail to capture nonlinear
hyperelasticity~{[}2{]},~{[}18{]} and high-fidelity soft-material
simulation is rarely available, robots must learn their own dynamics on
the fly~{[}11{]}, guided by uncertainty-aware active exploration to
acquire generalisable models efficiently~{[}12{]}.

\textbf{Proposed Approach. }Leveraging the Julia--Turing ecosystem, we
engineer a Bayesian adaptation framework employing differentiable
programming for joint state-inference and hyperparameter optimisation at
kilohertz frequencies, enabling real-time GP uncertainty propagation
across the MPC horizon.

\textbf{II. RELATED WORK}

GPs and neural networks are established nonparametric tools for
modelling nonlinear dynamics with uncertainty quantification.
Kocijan~{[}1{]},~{[}2{]}, Pillonetto et\,al.~{[}3{]},~{[}4{]}, and
Ljung~{[}5{]} provided the foundational theory for GP-based system
identification.

\emph{\textbf{A. Scalability and Online Learning}}

\begin{itemize}
\item
  \begin{quote}
  Sparse approximations: FITC (Snelson~\& Ghahramani) and VFE
  (Titsias~{[}6{]}) produce compact representations. Sparse Online GPs
  (Csátó~\& Opper) and streaming variants (Bui et\,al.) support
  continuous adaptation.
  \end{quote}
\item
  \begin{quote}
  Localised models: Local GP Regression (Nguyen-Tuong~\&
  Peters~{[}8{]},~{[}9{]}) partitions the input space incrementally for
  real-time learning.
  \end{quote}
\item
  \begin{quote}
  Recursive state-space models: RGPSSM (Zheng et\,al.~{[}5{]}) and
  related work~{[}7{]} adapt online to novel conditions via dynamic
  inducing-point sets.
  \end{quote}
\end{itemize}

\emph{\textbf{B. Applications in Control and Robotics}}

\begin{itemize}
\item
  \begin{quote}
  GP-MPC: Hewing et\,al.~{[}13{]} introduced chance-constrained GP-MPC
  with Taylor linearisation; Klenske et\,al.~{[}14{]} addressed periodic
  error correction.
  \end{quote}
\item
  \begin{quote}
  Safe RL~\& active exploration: Berkenkamp et\,al.~{[}19{]} pioneered
  Lyapunov-safe MBRL using GPs; SoftAE~{[}12{]} (Sukhija et\,al.)
  extends this to task-agnostic exploration for soft robots.
  \end{quote}
\end{itemize}

\emph{\textbf{C. Soft Robotics}}

\begin{itemize}
\item
  \begin{quote}
  Vision-based control: Fang et\,al.~{[}17{]} applied local GP
  regression for online kinematic control of soft arms.
  \end{quote}
\item
  \begin{quote}
  Inverse statics: Giorelli et\,al. addressed nonconstant-curvature
  manipulators via Jacobian neural methods.
  \end{quote}
\item
  \begin{quote}
  Sparse dynamics discovery: Fasel et\,al.~{[}18{]} used Ensemble-SINDy
  for robust equation extraction under high noise and limited data.
  \end{quote}
\end{itemize}

\textbf{III. SYSTEM ARCHITECTURE AND METHODOLOGY}

The framework is implemented in Julia~v1.11+ with the SciML ecosystem.
Nominal actuator physics are encoded as continuous ODEs in
PhysicsPriors.jl; unmodelled dynamics are captured by a Latent Force
Model (LFM) via Turing.jl. Exact gradients of the MPC objective with
respect to the control sequence are computed by Zygote.jl and
SciMLSensitivity.jl, backpropagating through the ODE solver without
finite-difference approximations.

\includegraphics[width=3.65486in,height=2.56111in]{media/image8.png}

\emph{Fig.~1. Closed-loop real-time adaptation architecture{[}39{]}.}

\emph{\textbf{A. Theoretical Foundations}}

\emph{1) Latent Force Models}

Alvarez, Luengo, and Lawrence~{[}20{]} embed GP priors inside
differential operators, providing a principled hybrid of physical
structure and data-driven residuals which is the mathematical core of
the VLFM.

Following the non-parametric formulation of Alvarez et al. [20], the non-stationary residual dynamics $f_{\text{latent}}(t)$ are driven by a zero-mean Gaussian Process (GP) coupled to a continuous temporal operator:
\begin{equation}
f_{\text{latent}}(t) \sim \mathcal{GP}\left(0, \kappa(t, t')\right)
\label{eq:latent_force_gp}
\end{equation}
To guarantee smooth, infinitely differentiable latent realizations, the covariance function $\kappa(t, t')$ is governed by a Squared Exponential (SE) kernel profile:
\begin{equation}
\kappa(t, t') = \sigma_{\text{gp}}^2 \exp\left(-\frac{(t - t')^2}{2\ell^2}\right)
\label{eq:se_kernel}
\end{equation}
parameterized by the signal variance $\sigma_{\text{gp}}^2$ and characteristic temporal lengthscale $\ell$.

\emph{2) Universal Differential Equations}

Rackauckas et\,al.~{[}21{]} establish that JIT-compiled adjoint
sensitivities in Julia outpace Python implementations for scientific ML
tasks, justifying our use of DifferentialEquations.jl with full-system
differentiability.

\emph{3) Physics-Informed Priors for Soft Robotics}

Thuruthel et\,al.~{[}2{]} demonstrate that soft-robot viscoelastic
properties defy analytical modelling --- the precise gap the
VLFM\textquotesingle s hybrid physics-data approach is designed to
bridge.

The continuous-time structural nominal dynamics of the soft actuator are explicitly parameterized as a non-linear Duffing-type mass-spring-damper system:
\begin{equation}
m\ddot{x}(t) + c\dot{x}(t) + kx(t) + 0.1kx^3(t) = u(t) + f_{\text{latent}}(t)
\label{eq:system_dynamics}
\end{equation}
where $m$, $c$, and $k$ denote the mass, viscoelastic damping, and nominal stiffness constants respectively, $u(t)$ is the continuous control input force, and $f_{\text{latent}}(t)$ represents the unmodeled morphological or structural deviations.
The predictive controller enforces a receding horizon control law optimized over $N$ discrete lookahead intervals of duration $\Delta t$. The deterministic optimization objective $J(u)$ targets aggressive position-tracking minimization while mitigating extreme control expenditure:
\begin{align}
\min_{u} \quad & J(u) = \sum_{i=1}^{N} \left( 10.0 \cdot (x_i - x_{\text{target}})^2 + 0.01 \cdot u_i^2 \right) \label{eq:mpc_objective} \\
\text{s.t.} \quad & u_i \in [-15.0, 15.0], \quad \forall i \in \{1, \dots, N\} \label{eq:mpc_constraints}
\end{align}
Gradients of the objective function w.r.t the control sequence ($\nabla_u J$) are backpropagated through a custom explicit Runge-Kutta 4th-order (RK4) scheme utilizing the automated differentiation (AD) source-to-source engine provided by Zygote.jl.

\emph{4) GP-MPC Baseline Context}

Hewing et\,al.~{[}13{]} provide a discrete-time, data-heavy GP-MPC
baseline. The VLFM differentiates itself through continuous-time,
physics-constrained modelling with superior sample efficiency.

\emph{\textbf{B. Competitor Baselines}}

\emph{Baseline~1 --- Sparse GP (Black-Box)}

Method: Variational Sparse GP (Titsias~{[}6{]}), ignoring ODE physics.
\emph{Expected result:} highest MSE; no generalisation to unseen state
regions due to sample inefficiency.

\emph{Baseline~2 --- Discrete GP-SSM}

Method: Transition model xₜ₊₁\,=\,f(xₜ)\,+\,ε with Euler integration
(Doerr et\,al.~{[}23{]}; Liu et\,al.~{[}7{]}).

\emph{Expected result:} compounding discretisation errors diverge over
long MPC horizons.

\emph{Baseline~3 --- Neural ODE}

Method: Neural network as ODE right-hand side (Chen et\,al.~{[}22{]}).
\emph{Expected result:} absence of epistemic uncertainty bounds causes
oscillations and overshoot in unexplored regions of the state space.

\textbf{IV. EXPERIMENTAL SETUP}

\emph{\textbf{A. Software Stack}}

The system architecture is compiled natively within the Julia v1.11+ environment [40] to exploit high-performance, JIT-compiled mathematical structures. Continuous-time initial value problems are resolved using DifferentialEquations.jl [41], with exact parametric sensitivities evaluated through SciMLSensitivity.jl and the Zygote.jl reverse-mode automatic differentiation compiler [42]. Nonparametric Bayesian tracking and online state estimation are executed via the Hamilton Monte Carlo (HMC) No-U-Turn Sampler (NUTS) within the Turing.jl probabilistic programming network [43].

\emph{\textbf{B. Baseline KNN Parameter Study}}

A K-Nearest Neighbours baseline was evaluated across neighbourhood sizes
to establish a data-driven reference. K\,=\,3 overfit localised sensor
noise, yielding suboptimal prediction. MSE was minimised at K\,=\,8,
adopted as the baseline hyperparameter.

\includegraphics[width=3.33333in,height=2.13889in]{media/image14.png}

\emph{Fig.~2 --- KNN Sensitivity Analysis Line chart: MSE vs.\,K
(K\,=\,1\ldots20) Vertical dashed marker at K\,=\,8}. \emph{KNN
hyperparameter study. Error minimised at K\,=\,8.}

\emph{\textbf{C. Closed-Loop Trajectory Tracking Protocol}}

The reference trajectory was defined as xₜₐᵣᵏᵉᵗ(t)\,=\,sin(1.5t),
simulating a rhythmic bending task. The actuator was initialised from
rest (x₀\,=\,0, v₀\,=\,0) over a 6.0\,s horizon with MPC timestep
Δt\,=\,0.1\,s and prediction horizon N\,=\,3. To stress-test adaptation,
all experiments used a deliberate prior mismatch: c\,=\,0.1,\,k\,=\,0.5
(true values: cₑₙᵤᵉ\,=\,0.7, kₑₙᵤᵉ\,=\,3.0). Control inputs were bounded
to ±15.0 and optimised over 50 gradient steps per timestep.

\textbf{Phase~1 --- Prior Exploration (0.0--3.0\,s): }The MPC operates
exclusively on the inaccurate prior. Control effort is systematically
insufficient to overcome the true actuator stiffness, establishing a
degraded-performance baseline.

\textbf{Phase~2 --- Posterior Adaptation (3.0--6.0\,s): }At
t\,=\,3.0\,s, the NUTS sampler assimilates Phase~1 sensor history and
updates the model posterior. The controller immediately integrates the
learned dynamics.

\textbf{V. RESULTS AND DISCUSSION}

\emph{\textbf{A. Platform Benchmarking}}

Table~I compares the proposed Julia implementation against the Python
ecosystem (GPyTorch/Pyro) across key performance dimensions.

\textbf{Table~I Platform Comparison}

\begin{longtable}[]{@{}
  >{\raggedright\arraybackslash}p{(\columnwidth - 4\tabcolsep) * \real{0.2625}}
  >{\raggedright\arraybackslash}p{(\columnwidth - 4\tabcolsep) * \real{0.3392}}
  >{\raggedright\arraybackslash}p{(\columnwidth - 4\tabcolsep) * \real{0.3982}}@{}}
\toprule\noalign{}
\begin{minipage}[b]{\linewidth}\raggedright
\textbf{Feature}
\end{minipage} & \begin{minipage}[b]{\linewidth}\raggedright
\textbf{Python (GPyTorch/Pyro)}
\end{minipage} & \begin{minipage}[b]{\linewidth}\raggedright
\textbf{Julia\,+\,Turing.jl}
\end{minipage} \\
\begin{minipage}[b]{\linewidth}\raggedright
Inference speed
\end{minipage} & \begin{minipage}[b]{\linewidth}\raggedright
GIL-limited; batch-oriented
\end{minipage} & \begin{minipage}[b]{\linewidth}\raggedright
Native JIT; streaming-capable
\end{minipage} \\
\begin{minipage}[b]{\linewidth}\raggedright
ODE integration
\end{minipage} & \begin{minipage}[b]{\linewidth}\raggedright
Requires custom bridging
\end{minipage} & \begin{minipage}[b]{\linewidth}\raggedright
Seamless via SciML
\end{minipage} \\
\begin{minipage}[b]{\linewidth}\raggedright
Real-time updates
\end{minipage} & \begin{minipage}[b]{\linewidth}\raggedright
Batch processing
\end{minipage} & \begin{minipage}[b]{\linewidth}\raggedright
Recursive, on-the-fly
\end{minipage} \\
\begin{minipage}[b]{\linewidth}\raggedright
Differentiation scope
\end{minipage} & \begin{minipage}[b]{\linewidth}\raggedright
Network layers only
\end{minipage} & \begin{minipage}[b]{\linewidth}\raggedright
Full system: model\,+\,solver\,+\,control
\end{minipage} \\
\midrule\noalign{}
\endhead
\bottomrule\noalign{}
\endlastfoot
\end{longtable}

\emph{\textbf{B. Execution Time Benchmarking}}

Forward-mode AD (Zygote.jl) through the Tsit5 ODE solver was benchmarked
against a discrete-time Euler approximation. Results are summarised in
Table~II.

\textbf{Table~II Execution Time and Memory}

\begin{longtable}[]{@{}
  >{\raggedright\arraybackslash}p{(\columnwidth - 6\tabcolsep) * \real{0.3000}}
  >{\raggedright\arraybackslash}p{(\columnwidth - 6\tabcolsep) * \real{0.1969}}
  >{\raggedright\arraybackslash}p{(\columnwidth - 6\tabcolsep) * \real{0.1969}}
  >{\raggedright\arraybackslash}p{(\columnwidth - 6\tabcolsep) * \real{0.3062}}@{}}
\toprule\noalign{}
\begin{minipage}[b]{\linewidth}\raggedright
\textbf{Method}
\end{minipage} & \begin{minipage}[b]{\linewidth}\raggedright
\textbf{Time}
\end{minipage} & \begin{minipage}[b]{\linewidth}\raggedright
\textbf{Memory}
\end{minipage} & \begin{minipage}[b]{\linewidth}\raggedright
\textbf{Stability}
\end{minipage} \\
\begin{minipage}[b]{\linewidth}\raggedright
Euler (discrete)
\end{minipage} & \begin{minipage}[b]{\linewidth}\raggedright
17.92\,μs
\end{minipage} & \begin{minipage}[b]{\linewidth}\raggedright
13.6\,KiB
\end{minipage} & \begin{minipage}[b]{\linewidth}\raggedright
Low (Δt\,→\,0)
\end{minipage} \\
\begin{minipage}[b]{\linewidth}\raggedright
SciML / Proposed
\end{minipage} & \begin{minipage}[b]{\linewidth}\raggedright
2.06\,ms
\end{minipage} & \begin{minipage}[b]{\linewidth}\raggedright
557\,KiB
\end{minipage} & \begin{minipage}[b]{\linewidth}\raggedright
High (unconditional)
\end{minipage} \\
\midrule\noalign{}
\endhead
\bottomrule\noalign{}
\endlastfoot
\end{longtable}

Although the Euler step is faster per evaluation (17.92\,μs), matching
the continuous solver\textquotesingle s numerical stability for
viscoelastic dynamics would require reducing Δt by two orders of
magnitude, eliminating the speed advantage. The proposed SciML approach
operates at ∼ 480\,Hz within a 600\,KiB memory footprint, confirming
viability for embedded deployment.

\emph{\textbf{C. Bayesian Parameter Estimation (NUTS)}}

Online estimation via NUTS successfully converged to the true physical
parameters (damping\,c, stiffness\,k), effectively decoupling them from
unmodelled disturbances captured by the GP. Trace plots confirmed
healthy Markov chain mixing, demonstrating reliable uncertainty
quantification without catastrophic forgetting which is a known failure
mode of discrete-time neural network approaches.

\includegraphics[width=3.33333in,height=2.08333in]{media/image20.png}

\includegraphics[width=3.33333in,height=3.08854in]{media/image7.png}

\emph{Fig.~3 --- Posterior Densities See posterior\_distributions.png in
the code repository {[}39{]} (a)\,Prior vs.\,posterior for damping c;
(b)\,prior vs.\,posterior for stiffness k. Vertical dashed lines
indicate true parameter values. Trace plots confirm HMC mixing.}
\emph{NUTS posterior distributions for c (damping) and k (stiffness).}

\emph{\textbf{D. Closed-Loop Trajectory Tracking}}

\emph{Phase~1 (0--3\,s) --- Degraded Performance}

With the inaccurate prior (k\,=\,0.5 vs.\,kₑₙᵤᵉ\,=\,3.0), the MPC
generates insufficient control effort, producing a persistent phase lag
in the tracked trajectory.

\emph{Trigger (t\,=\,3\,s) --- Bayesian Update}

The NUTS pipeline assimilates accumulated sensor history and rapidly
converges on the true stiffness and damping. The MPC posterior is
updated instantaneously.

\emph{Phase~2 (3--6\,s) --- Tight Tracking}

The phase lag vanishes immediately post-update. The state trajectory
hugs the reference, confirming that continuous-time SciML gradients
coupled with Bayesian inference enable robust real-time adaptation.

\includegraphics[width=3.33333in,height=2.08333in]{media/image15.png}

\emph{Fig.~4 --- Trajectory Tracking See baseline\_comparison.png in the
code repository {[}39{]} Dashed black = reference; solid blue = robot
state; red vertical = t\,=\,3\,s Shaded regions: Phase~1 (light red),
Phase~2 (light green). Inset: transition detail.Closed-loop tracking
with Bayesian adaptation at t\,=\,3\,s}

\emph{\textbf{E. Regularisation and Informative Priors}}

Setting the GP lengthscale prior to ℓ\,∼\,LogNormal(0,\,0.25)
regularised sensitivity to high-frequency sensor noise, yielding a
31\,\% MSE improvement over the unregularised baseline while preserving
sensitivity to genuine morphological shifts.

\emph{\textbf{F. Comparative Baseline Analysis}}

\textbf{Table~III summarises trajectory-tracking MSE across all
evaluated methods.}

\begin{longtable}[]{@{}
  >{\raggedright\arraybackslash}p{(\columnwidth - 4\tabcolsep) * \real{0.3963}}
  >{\raggedright\arraybackslash}p{(\columnwidth - 4\tabcolsep) * \real{0.1921}}
  >{\raggedright\arraybackslash}p{(\columnwidth - 4\tabcolsep) * \real{0.4116}}@{}}
\toprule\noalign{}
\begin{minipage}[b]{\linewidth}\raggedright
\textbf{Method}
\end{minipage} & \begin{minipage}[b]{\linewidth}\raggedright
\textbf{MSE}
\end{minipage} & \begin{minipage}[b]{\linewidth}\raggedright
\textbf{Key Limitation}
\end{minipage} \\
\begin{minipage}[b]{\linewidth}\raggedright
Sparse GP (black-box)
\end{minipage} & \begin{minipage}[b]{\linewidth}\raggedright
0.1302
\end{minipage} & \begin{minipage}[b]{\linewidth}\raggedright
No physics prior; sample-inefficient
\end{minipage} \\
\begin{minipage}[b]{\linewidth}\raggedright
Discrete GP-SSM
\end{minipage} & \begin{minipage}[b]{\linewidth}\raggedright
0.7790
\end{minipage} & \begin{minipage}[b]{\linewidth}\raggedright
Compounding discretisation errors
\end{minipage} \\
\begin{minipage}[b]{\linewidth}\raggedright
Neural ODE
\end{minipage} & \begin{minipage}[b]{\linewidth}\raggedright
0.0155
\end{minipage} & \begin{minipage}[b]{\linewidth}\raggedright
No epistemic uncertainty; oscillations
\end{minipage} \\
\begin{minipage}[b]{\linewidth}\raggedright
VLFM (proposed)
\end{minipage} & \begin{minipage}[b]{\linewidth}\raggedright
0.0107
\end{minipage} & \begin{minipage}[b]{\linewidth}\raggedright
--- best result
\end{minipage} \\
\midrule\noalign{}
\endhead
\bottomrule\noalign{}
\endlastfoot
\end{longtable}

\textbf{Table~III Trajectory-Tracking MSE Comparison}

\textbf{Summary. }The VLFM achieves a 92\,\% MSE reduction over the
sparse GP and a 31\,\% reduction over the Neural ODE by combining
continuous-time ODE stability, GP uncertainty quantification, and a
physics-informed prior over the residual latent force.

\textbf{VI. CODE AVAILABILITY}

The full implementation including the VLFM core (src/), execution-time
benchmarks (benchmarks/), training and evaluation scripts (scripts/),
and unit tests (test/) are publicly available under an open-source
license:

\href{https://github.com/rainerrodrigues/VLFM_SoftRobotics}{\ul{https://github.com/rainerrodrigues/VLFM\_SoftRobotics}}

The repository also contains the three result figures cited in this
paper (posterior\_distributions.png, baseline\_comparison.png,
closed\_loop\_hero\_graphic.png). A Julia environment file
(Project.toml) is provided for fully reproducible dependency management.

\textbf{VII. CONCLUSION AND FUTURE WORK}

\emph{\textbf{A. Summary}}

This paper presented a continuous-time Variational Latent Force Model
(VLFM) for real-time adaptive control of soft and deformable robots,
implemented in Julia with SciML and Turing.jl. A closed-loop experiment
under deliberate model misspecification validated the framework: the
VLFM assimilated streaming sensor data at t\,=\,3.0\,s via NUTS
sampling, converged on the true physical parameters, and immediately
reduced a severe phase-lag error to tight reference following. Operating
at ∼ 480\,Hz within a 600\,KiB memory footprint, the VLFM achieved the
lowest MSE (0.0107) among all evaluated methods, confirming suitability
for embedded deployment.

\emph{\textbf{B. Future Work}}

\begin{enumerate}
\def\labelenumi{\arabic{enumi}.}
\item
  \begin{quote}
  Extension to multi-segment soft robots under distributed sensing,
  testing RGPSSM scalability in higher-dimensional state spaces.
  \end{quote}
\item
  \begin{quote}
  Integration of distributional-shift detection to identify mechanical
  failures or environmental contacts, triggering targeted model
  re-identification.
  \end{quote}
\item
  \begin{quote}
  Hardware validation on pneumatic or cable-driven actuators, accounting
  for pressure delays, hysteresis, and thermal drift.
  \end{quote}
\item
  \begin{quote}
  Multi-fidelity learning combining cheap simulation data with sparse
  real-world observations to reduce sample complexity.
  \end{quote}
\item
  \begin{quote}
  Safe RL integration using GP uncertainty estimates as Lyapunov
  stability certificates for online policy improvement~{[}19{]}.
  \end{quote}
\end{enumerate}

\textbf{VIII. REFERENCES}

\textbf{{[}1{]} }J.\,Kocijan, A.\,Girard, and B.\,Banko, ``Dynamic
systems identification with Gaussian processes,'' Math. Comput. Model.
Dyn. Syst., vol.\,11, no.\,4, pp.\,411--424, 2005,
doi:\,10.1080/13873950500068567.

\textbf{{[}2{]} }J.\,Kocijan, Modelling and Control of Dynamic Systems
Using Gaussian Process Models, 1st\,ed. Cham: Springer, 2016,
doi:\,10.1007/978-3-319-21021-6.

\textbf{{[}3{]} }G.\,Pillonetto, F.\,Dinuzzo, T.\,Chen, and
G.\,De\,Nicolao, ``Kernel methods in system identification, machine
learning and function estimation: a survey,'' Automatica, vol.\,50,
no.\,3, pp.\,657--682, 2014, doi:\,10.1016/j.automatica.2014.01.001.

\textbf{{[}4{]} }A.\,Scampicchio, E.\,Arcari, A.\,Lahr, and
M.\,N.\,Zeilinger, ``Gaussian processes for dynamics learning in model
predictive control,'' Annu. Rev. Control, vol.\,60, p.\,101034, 2025,
doi:\,10.1016/j.arcontrol.2025.101034.

\textbf{{[}5{]} }T.\,Zheng, L.\,Cheng, S.\,Gong, and X.\,Huang,
``Recursive Gaussian process state space model,'' arXiv:2411.14679,
Nov.\,2024.

\textbf{{[}6{]} }M.\,Titsias, ``Variational learning of inducing
variables in sparse Gaussian processes,'' in Proc. AISTATS, 2009,
pp.\,567--574.

\textbf{{[}7{]} }Y.\,Liu and P.\,M.\,Djurić, ``Gaussian process
state-space models with time-varying parameters and inducing points,''
in Proc. EUSIPCO, 2021, pp.\,1462--1466.

\textbf{{[}8{]} }D.\,Nguyen-Tuong and J.\,Peters, ``Model learning for
robot control: a survey,'' Cogn. Process., vol.\,12, no.\,4,
pp.\,319--340, 2011, doi:\,10.1007/s10339-011-0404-1.

\textbf{{[}9{]} }D.\,Nguyen-Tuong, ``Model learning in robot control,''
Ph.D.\,dissertation, Max-Planck Inst. Biol. Cybernetics, Tübingen,
Germany, 2011.

\textbf{{[}10{]} }J.\,Lindinger et\,al., ``Free-form variational
inference for Gaussian process state-space models,'' arXiv preprint,
2023.

\textbf{{[}11{]} }G.\,Pillonetto and L.\,Ljung, ``Full Bayesian
identification of linear dynamic systems using stable kernels,'' Proc.
Natl. Acad. Sci. USA, vol.\,120, no.\,18, p.\,e2218197120, 2023,
doi:\,10.1073/pnas.2218197120.

\textbf{{[}12{]} }H.\,Zheng, B.\,Sukhija, C.\,Li, K.\,Iten, A.\,Krause,
and R.\,K.\,Katzschmann, ``Learning soft robotic dynamics with active
exploration,'' arXiv:2510.27428, Oct.\,2025.

\textbf{{[}13{]} }L.\,Hewing, J.\,Kabzan, and M.\,N.\,Zeilinger,
``Cautious model predictive control using Gaussian process regression,''
IEEE Trans. Control Syst. Technol., vol.\,28, no.\,6, pp.\,2736--2743,
2020, doi:\,10.1109/TCST.2019.2949757.

\textbf{{[}14{]} }E.\,D.\,Klenske, M.\,N.\,Zeilinger, B.\,Schölkopf, and
P.\,Hennig, ``Gaussian process-based predictive control for periodic
error correction,'' IEEE Trans. Control Syst. Technol., vol.\,24,
no.\,1, pp.\,110--121, 2016, doi:\,10.1109/TCST.2015.2420629.

\textbf{{[}15{]} }G.\,Williams, N.\,Wagener, B.\,Goldfain, P.\,Drews,
J.\,M.\,Rehg, B.\,Boots, and E.\,A.\,Theodorou, ``Information theoretic
MPC for model-based reinforcement learning,'' in Proc. ICRA, 2017,
pp.\,1714--1721, doi:\,10.1109/ICRA.2017.7989202.

\textbf{{[}16{]} }K.-H.\,Lee, D.\,K.\,C.\,Fu, C.\,W.\,M.\,Leong,
M.\,Chow, H.-C.\,Fu, K.\,Althoefer, K.\,Y.\,Sze, C.-K.\,Yeung, and
K.-W.\,Kwok, ``Nonparametric online learning control for soft continuum
robot: an enabling technique for effective endoscopic navigation,'' Soft
Robot., vol.\,4, no.\,4, pp.\,324--337, 2017,
doi:\,10.1089/soro.2016.0065.

\textbf{{[}17{]} }G.\,Fang, X.\,Wang, K.\,Wang, and K.-H.\,Lee,
``Vision-based online learning kinematic control for soft robots using
local Gaussian process regression,'' IEEE Robot. Autom. Lett., vol.\,4,
no.\,2, pp.\,1194--1201, 2019, doi:\,10.1109/LRA.2019.2893691.

\textbf{{[}18{]} }U.\,Fasel, J.\,N.\,Kutz, B.\,W.\,Brunton, and
S.\,L.\,Brunton, ``Ensemble-SINDy: robust sparse model discovery in the
low-data, high-noise limit,'' Proc. R. Soc. A, vol.\,478, no.\,2260,
p.\,20210904, 2022.

\textbf{{[}19{]} }F.\,Berkenkamp, M.\,Turchetta, A.\,P.\,Schoellig, and
A.\,Krause, ``Safe model-based reinforcement learning with stability
guarantees,'' in Proc. NeurIPS, vol.\,30, 2017.

\textbf{{[}20{]} }M.\,A.\,Alvarez, D.\,Luengo, and N.\,D.\,Lawrence,
``Latent force models,'' in Proc. AISTATS, 2009, pp.\,9--16.

\textbf{{[}21{]} }C.\,Rackauckas et\,al., ``Universal differential
equations for scientific machine learning,'' arXiv:2001.04385, 2020.

\textbf{{[}22{]} }R.\,T.\,Q.\,Chen, Y.\,Rubanova, J.\,Bettencourt, and
D.\,Duvenaud, ``Neural ordinary differential equations,'' in Proc.
NeurIPS, 2018.

\textbf{{[}23{]} }A.\,Doerr, C.\,Daniel, M.\,Schiegg, N.\,Deisenroth,
S.\,Schaal, M.\,Toussaint, and T.\,Rasch, ``Probabilistic recurrent
state-space models,'' in Proc. ICML, 2018.

\textbf{{[}24{]} }B.\,Ninness and S.\,Henriksen, ``Bayesian system
identification via Markov chain Monte Carlo techniques,'' Automatica,
vol.\,46, no.\,1, pp.\,40--51, 2010,
doi:\,10.1016/j.automatica.2009.10.015.

\textbf{{[}25{]} }L.\,Treven, P.\,Wenk, F.\,Dörfler, and A.\,Krause,
``Distributional gradient matching for learning uncertain neural
dynamics models,'' in Proc. NeurIPS, vol.\,34, 2021, pp.\,2719--2727.

\textbf{{[}26{]} }P.\,Wenk, A.\,Gotovos, S.\,Bauer, N.\,S.\,Gorbach,
A.\,Krause, and J.\,M.\,Buhmann, ``Fast Gaussian process based gradient
matching for parameter identification in systems of nonlinear ODEs,'' in
Proc. AISTATS, 2019, pp.\,1351--1360.

\textbf{{[}27{]} }M.\,Dubied, A.\,Lahr, M.\,N.\,Zeilinger, and
J.\,Köhler, ``A robust and adaptive MPC formulation for Gaussian process
models,'' arXiv:2507.02098, 2025.

\textbf{{[}28{]} }M.\,Greeff, A.\,W.\,Hall, and A.\,P.\,Schoellig,
``Learning a stability filter for uncertain differentially flat systems
using Gaussian processes,'' in Proc. CDC, 2021, pp.\,789--794.

\textbf{{[}29{]} }Y.\,Liu, P.\,Wang, and R.\,Tóth, ``Learning for
predictive control: a dual Gaussian process approach,'' Automatica,
vol.\,177, p.\,112316, 2025, doi:\,10.1016/j.automatica.2025.112316.

\textbf{{[}30{]} }A.\,Lederer, ``Gaussian processes in control:
performance guarantees through efficient learning,''
Ph.D.\,dissertation, TU München, Munich, Germany, 2023.

\textbf{{[}31{]} }S.\,Formentin, M.\,Mazzoleni, M.\,Scandella, and
F.\,Previdi, ``Nonlinear system identification via data augmentation,''
Syst. Control Lett., vol.\,128, pp.\,44--53, 2019,
doi:\,10.1016/j.sysconle.2019.04.004.

\textbf{{[}32{]} }J.\,Hendriks, J.\,R.\,Z.\,Holdsworth, A.\,Wills, and
T.\,B.\,Schön, ``Data to controller for nonlinear systems: an
approximate solution,'' arXiv:2103.08782, 2021.

\textbf{{[}33{]} }W.\,Cao and G.\,Pillonetto, ``Dealing with
collinearity in large-scale linear system identification using Bayesian
regularization,'' arXiv:2203.13633, 2022.

\textbf{{[}34{]} }M.\,Šmíd et\,al., ``Online learning and control for
data-augmented quadrotor model,'' arXiv:2304.00503, 2024.

\textbf{{[}35{]} }R.\,Du, R.\,Lin, Y.\,Shen, and M.\,Egerstedt, ``Online
learning and coverage of unknown fields using random-feature Gaussian
processes,'' arXiv:2509.08117, 2025.

\textbf{{[}36{]} }A.\,Romero, E.\,Aljalbout, Y.\,Song, and
D.\,Scaramuzza, ``Actor-critic model predictive control: differentiable
optimization meets reinforcement learning for agile flight,'' Robotics
and Perception Group, ETH Zürich, 2024.

\textbf{{[}37{]} }G.\,Williams, P.\,Drews, B.\,Goldfain, J.\,M.\,Rehg,
and E.\,A.\,Theodorou, ``Information-theoretic model predictive control:
theory and applications to autonomous driving,'' IEEE Trans. Robot.,
vol.\,34, no.\,6, pp.\,1603--1622, Dec.\,2018.

\textbf{{[}38{]} }P.\,L.\,Green, L.\,J.\,Devlin, R.\,E.\,Moore, and
R.\,J.\,Jackson, ``Increasing the efficiency of sequential Monte Carlo
samplers through approximately optimal L-kernels,'' Mech. Syst. Signal
Process., vol.\,162, p.\,108028, 2022,
doi:\,10.1016/j.ymssp.2021.108028.

\textbf{{[}39{]} }R.\,Rodrigues, ``VLFM\_SoftRobotics,'' GitHub
repository, 2025. {[}Online{]}. Available:
\ul{https://github.com/rainerrodrigues/VLFM\_SoftRobotics}

{[}40{]} J. Bezanson, A. Edelman, S. Karpinski, and V. B. Shah, ``Julia:
A fresh approach to numerical computing,'' SIAM Review, vol. 59, no. 1,
pp. 65--98, 2017.

{[}41{]} C. Rackauckas and Q. Nie, ``Differentialequations.jl---a
performant and feature-rich ecosystem for solving differential equations
in Julia,'' Journal of Open Source Software, vol. 2, no. 15, p. 15,
2017.

{[}42{]} M. Innes, ``Don\textquotesingle t unroll adjoint:
Differentiating SSA-form programs,'' in Proc. Advances in Neural
Information Processing Systems (NeurIPS), 2018, pp. 7994--8004.

{[}43{]} H. Ge, K. Xu, and Z. Ghahramani, ``Turing: a language for
flexible probabilistic inference,'' in Proc. International Conference on
Artificial Intelligence and Statistics (AISTATS), 2018, pp. 1682--1690.

\end{document}
